\subsection{Gamification}
Studien innefattar inte någon praktisk instans av gamification utan reflekterar snarare kring dess tillämpning i en C-DAS-OB-applikation och vad som är viktigt kopplat till gamification ur lokförarens perspektiv.
Detta i enlighet med studiens syfte, omfattning och forskningsfrågor.
%När gamification ska tillämpas i C-DAS-applikationer är det viktigt att det inte görs på ett sätt så att lokförarna inte känner sig utvärderade utan istället med hjälp av gamification kan utvärdera sig själva.
Tillämpandet av gamification och intrinsisk motivation utanför spelvärlden är en relativt ny trend och sett till det förvånansvärt passionerade engagemang dataspel frambringar kan potentialen mycket väl vara stor vid rätt utförande.
Gamification är dock komplext och kräver djuplodande kunskaper i såväl beteendevetenskap som datavetenskap och inte minst domänkunskap för slutanvändare och det system varuti det är tänkt att implementeras.
\paragraph{}
\hspace{-11pt}Även om C-DAS med gamification-attribut ger god möjlighet för management att följa upp lokförarnas efterlevnad av körprofiler finns det risk att detta får en stor negativ inverkan på arbetsmiljön om förare känner sig utvärderade eller övervakade~\cite{Bae2014}.
Det är bättre om management tar ett steg tillbaka beträffande feedback för att istället låta medarbetaren själv utvärdera sig och hitta förbättringsmöjligheter~\cite{Cwil2018}.
Det gäller i synnerhet negativ feedback men även positiv dito då all form av uppföljning skulle kunna leda till känslan av att vara övervakad.
Påståendet styrk både i litteraturstudien och i uttalanden från flera respondenter i denna studies intervjuer.
Däremot kan andra potentiella nyckelfaktorer skönjas ur ett gamificationperspektiv kopplat till DAS-applikationer.
Exempelvis självreflektion och möjligheten att utvärdera egna handlingar kopplat till tågets framförande såsom punktlighet och energiförbrukning samt jämförelsen med andra lokförare eller handlingar vid en viss situation~\cite{Bae2014, Cwil2018}.
Det kan exempelvis röra sig om att skapa egna \lq{}grupper\rq{} där medlemmar tävlar individuellt eller i grupp om högst följsamhet av körprofil, bäst punktlighet eller lägst energikonsumtion~\cite{Bae2014}.
Att analysera förarnas agerande i olika situationer kan också vara bra underlag för utbildningsmaterial.
Det är som tidigare nämnt viktigt att management inte använder statistiken för att följa upp eller mäta lokförarnas färdigheter.
Det är också viktigt att eventuell följsamhetsspårning är korrekt implementerad och mäter efterlevnad av körprofil mycket precist och med högt förtroende.
%Tågläge styrker och legitimerar körprofil
\paragraph{}
\hspace{-11pt}Bör inte vara jobbigt.
Kan implementeras stegvis om flera steg föreligger.
